\chapter{Hypersomnia}

\section*{Definition}
Hypersomnia is characterized by excessive daytime sleepiness (EDS) or prolonged nighttime sleep despite adequate sleep duration, with an Epworth Sleepiness Scale score typically \textgreater{}10 or a mean sleep latency of \textless{}8 minutes on the Multiple Sleep Latency Test.

\section*{What Happens in Your Body}
Dysfunction of hypothalamic hypocretin/orexin signaling, dysregulation of GABAergic and histaminergic arousal pathways, or impaired circadian rhythm entrainment leads to inability to maintain wakefulness. In narcolepsy type~1, autoimmune destruction of hypocretin-producing neurons in the lateral hypothalamus is the primary mechanism.

\section*{Causes}
\begin{itemize}
  \item \textbf{Primary sleep disorders:} Narcolepsy (type 1 and 2), idiopathic hypersomnia, Kleine--Levin syndrome
  \item \textbf{Secondary medical causes:} Obstructive sleep apnea, circadian rhythm disorders, Parkinson disease, multiple sclerosis, brain tumors, head trauma
  \item \textbf{Psychiatric:} Atypical depression, bipolar disorder, seasonal affective disorder
  \item \textbf{Drug-induced:} Sedative--hypnotics, antihistamines, antidepressants (trazodone, mirtazapine), antipsychotics, alcohol, cannabis
  \item \textbf{Nutritional/metabolic:} Hypothyroidism, iron deficiency, B\(\_{12}\) deficiency, obesity--hypoventilation syndrome
\end{itemize}

\section*{When to See a Doctor}
See your doctor if:
\begin{itemize}
  \item Falling asleep while driving, operating machinery, or during conversations
  \item Sleep attacks without warning or cataplexy (sudden bilateral muscle weakness triggered by emotion)
  \item Hypersomnia persists despite \textgreater{}8 hours of nocturnal sleep and optimized sleep hygiene
  \item Bed partner reports loud snoring, gasping, or witnessed apneas
\end{itemize}

\section*{Related Symptoms}
\begin{itemize}
  \item Cataplexy, sleep paralysis, hypnagogic/hypnopompic hallucinations, fatigue, brain fog, cognitive slowing, headache upon awakening, restless legs symptoms
\end{itemize}
\textbf{Important:} Hypersomnia must be distinguished from simple fatigue or lack of motivation, as the management of primary hypersomnia disorders differs substantially from secondary causes.

\section*{Self-Care / Home Management}
\begin{itemize}
  \item Maintain a consistent sleep--wake schedule 7 days per week, including weekends
  \item Incorporate 1--2 strategic naps of 15--20 minutes in the early afternoon
  \item Avoid alcohol, caffeine after 2 PM, and heavy meals before bed
  \item Use bright-light exposure (\textgreater{}2,500~lux) for 30 minutes upon awakening
  \item Remove electronic screens and minimize noise/light in the bedroom
\end{itemize}

\section*{How Your Doctor Will Evaluate You}
\begin{itemize}
  \item Overnight polysomnography to rule out sleep-disordered breathing and periodic limb movements
  \item Multiple Sleep Latency Test (MSLT) to quantify sleep propensity and detect sleep-onset REM periods
  \item Actigraphy for 1--2 weeks to evaluate rest--activity patterns and circadian alignment
  \item Laboratory evaluation: TSH, complete blood count, ferritin, vitamin B\(\_{12}\), HIV/hepatitis serologies if indicated
  \item Brain MRI with pituitary protocol if central hypersomnia of unclear causes
\end{itemize}

\section*{Treatment Options}
\begin{itemize}
  \item \textbf{Pharmacologic:} Modafinil or armodafinil as first-line stimulants; sodium oxybate for narcolepsy with cataplexy; methylphenidate or amphetamine salts as second-line agents; pitolisant (histamine H\(\_3\)-receptor inverse agonist) for narcolepsy
  \item \textbf{Non-pharmacologic:} Scheduled naps, cognitive behavioral therapy for hypersomnia (CBT-H), CPAP if OSA is identified, treatment of underlying endocrinopathy or deficiency
  \item \textbf{Lifestyle:} Driving restrictions until symptom control is achieved; referral to sleep medicine specialist; patient education about medication side effects and tolerance
\end{itemize}

\section*{References}

\begin{thebibliography}{99}
\bibitem{hypersomnia:icsd3} American Academy of Sleep Medicine. \emph{International Classification of Sleep Disorders}. 3rd ed, Text Revision. AASM; 2023. Central Disorders of Hypersomnolence.
\bibitem{hypersomnia:harrison-sleep} Longo DL, et al. \emph{Harrison's Principles of Internal Medicine}. 21st ed. McGraw-Hill; 2022. Chapter 28: Hypersomnia.
\bibitem{hypersomnia:uptodate-hyper-sleep} Scammell TE, et al. Approach to the patient with hypersomnia. In: UpToDate, Post TW (Ed), UpToDate, Waltham, MA; 2025.
\bibitem{hypersomnia:bassetti} Bassetti CLA, et al. Narcolepsy and idiopathic hypersomnia. \emph{Lancet Neurol}. 2023;22(7):611-625.
\bibitem{hypersomnia:kryger} Kryger MH, et al. \emph{Principles and Practice of Sleep Medicine}. 7th ed. Elsevier; 2022. Chapter 87: Idiopathic Hypersomnia.
\bibitem{hypersomnia:trotti} Trotti LM, et al. Idiopathic hypersomnia: diagnosis and management. \emph{Chest}. 2023;163(3):680-689.
\end{thebibliography}
